\documentclass[11pt]{article}
\usepackage{amsmath}
\usepackage{amssymb}
\usepackage{graphicx}
\usepackage{geometry}
\geometry{a4paper, margin=1in}

\title{Stage 2 Interim Report: Sparse-Sensor Reconstruction of Aerodynamic Flow Fields Using Reduced-Order Modeling and Deep Neural Networks}
\author{VRIGHT BROTHERS: Vedant Lohare (241151) \& Rajvardhan Beniwal (240838)\\
Department of Aerospace Engineering, IIT Kanpur\\
AE646: Scientific Machine Learning for Fluid Mechanics}
\date{September 2026}

\begin{document}

\maketitle

\section{Introduction and Objectives}
This report details the Stage 2 implementation for the sparse-sensor reconstruction of high-fidelity aerodynamic flow fields. The primary objective is to solve the under-determined spatial inverse problem:
\begin{equation}
\mathbf{y} = \mathbf{C}\mathbf{x} + \boldsymbol{\eta}, \quad \mathbf{x} \in \mathbb{R}^N, \quad \mathbf{y} \in \mathbb{R}^p, \quad \boldsymbol{\eta} \sim \mathcal{N}(0, \sigma^2 \mathbf{I})
\end{equation}
where $p \ll N$. We benchmark Tikhonov-regularized Gappy POD with Q-DEIM against a custom Physics-Aware Sensor-to-Field MLP.

\section{Methodology}
\subsection{Classical Baseline: Gappy POD with Q-DEIM}
We construct a mean-centered fluctuation snapshot matrix $\tilde{\mathbf{X}} \in \mathbb{R}^{N \times M}$ and compute its thin SVD to extract dominant spatial modes $\boldsymbol{\Phi}_r$. When $p < r$, the point-collocation matrix $\mathbf{M} = \boldsymbol{\Phi}_r^T \mathbf{C}^T \mathbf{C} \boldsymbol{\Phi}_r$ is strictly rank-deficient. Thus, modal coefficients $\mathbf{a}$ are recovered by minimizing the regularized spatial residual:
\begin{equation}
\hat{\mathbf{a}} = \left( \boldsymbol{\Phi}_r^T \mathbf{C}^T \mathbf{C} \boldsymbol{\Phi}_r + \mu \mathbf{I}_r \right)^{-1} \boldsymbol{\Phi}_r^T \mathbf{C}^T (\mathbf{y} - \mathbf{C}\bar{\mathbf{x}})
\end{equation}
Optimal sensor placement is determined using Q-DEIM, which performs a pivoted QR decomposition on $\boldsymbol{\Phi}_r^T$.

\subsection{Deep Learning Architecture: SensorMLP}
We train a 5-layer feedforward network $f_\theta : \mathbb{R}^p \rightarrow \mathbb{R}^N$ parametrized by intermediate latent dimensions $[128 \rightarrow 256 \rightarrow 512 \rightarrow 1024 \rightarrow N]$ with LayerNorm and GELU activations. The network is regularized using a composite physics-aware loss:
\begin{equation}
\mathcal{L}_{\text{total}} = \text{MSE}(\mathbf{\hat{x}}, \mathbf{x}) + \gamma_1 \|\nabla_x \mathbf{\hat{x}} - \nabla_x \mathbf{x}\|_1 + \gamma_2 \|\nabla_y \mathbf{\hat{x}} - \nabla_y \mathbf{x}\|_1
\end{equation}

\section{Implementation Status}
We have successfully constructed the \texttt{data/}, \texttt{src/rom/}, \texttt{src/sciml/}, and \texttt{src/evaluation/} Python modules. The PyTorch training pipeline handles the composite gradients and backpropagates through the spatial finite-differences efficiently.

\section{AI-Tool Use Declaration}
Large language models (Gemini) were consulted for code modularization, repository structuring, and LaTeX syntax formatting. All mathematical derivations, implementations, and physical interpretations are fully understood and owned by the authors.

\end{document}
